\makeatletter
\def\input@path{{./}{../../../preamble/}{book2026-2027/preamble/}}
\makeatother
\input{preamble.tex}
\usepackage{needspace}
\hypersetup{pdftitle={Grade 10 Term 2 Discrete Probability}}
\setlength{\emergencystretch}{2em}
\title{Grade 10, Term 2, Lesson 1, Criterion C3\\Bernoulli Indicators and Financial Consequences}
\author{}
\date{}
\begin{document}
\maketitle
\setcounter{tocdepth}{1}
\setcounter{secnumdepth}{2}
\phantomsection\label{toc}
\tableofcontents
\section*{Evaluated criterion}
\textbf{C3:} Models binary economic outcomes with Bernoulli random variables, interpreting success and failure in contexts such as default, gain, loss, or claim occurrence.
\section{Introduction}\label{sec:t2-c3-introduction}
All data in this activity are hypothetical teaching models. Probabilities have no monetary units. Keep exact expressions where useful and round final probabilities to four decimal places.
\subsection{Label one uncertain event}
A school enterprise sends one invoice. Define $X=1$ if it is paid on time and $X=0$ otherwise. This is an \emph{indicator}: its number records whether a precisely defined event happens. A Bernoulli variable describes one binary observation. Success is the event labelled 1; failure is its complement, labelled 0.
\textbf{Check.} Must success be financially good? \emph{Suggested answer:} No. Default can be labelled success mathematically.


\subsection{Probability and support}
Write $p=P(X=1)$, read as p is the probability that X equals one. Then $P(X=0)=1-p$. These probabilities have no units and add to 1. For $0<p<1$ the support is $\{0,1\}$. If $p=0$ or 1, only one label has positive probability.
\textbf{Check.} If $p=0.08$, what is $P(X=0)$? \emph{Suggested answer:} $1-0.08=0.92$.


\subsection{Separate the label from the consequence}
If an insured event costs 200 USD, $L=200X$ USD measures its loss. The indicator is unitless: $X=1$ records occurrence, while $L=200$ measures money. Reversing success gives $Z=1-X$ with $P(Z=1)=1-p$. A single observation needs no independence assumption; independence matters when combining several observations.
\textbf{Check.} If $X=1$ means default, what does $X=0$ mean? \emph{Suggested answer:} No default during the stated observation period.

\textbf{Check.} Is a 200 USD loss indicator equal to 200? \emph{Suggested answer:} No. The indicator is 1; its separate monetary consequence is 200 USD.

\section{General worked problems}
\Needspace{14\baselineskip}
\subsection{Introductory general problem: one binary trial}
An event A has probability 0.3. Define $X=1$ if A occurs and 0 otherwise.

\textbf{Questions/Tasks}
\begin{enumerate}
\item Give the support and calculate both masses from the Bernoulli formula.
\end{enumerate}
\subsubsection*{Complete worked solution}
\begin{center}
\begin{tabular}{llp{0.48\linewidth}}
\toprule
Symbol & Value & Meaning in this problem \\ \midrule
$p$ & $0.3$ & Probability of event A \\
$X$ & $0$ or $1$ & Indicator value; no monetary units \\
\bottomrule
\end{tabular}
\end{center}

\par\noindent\begin{minipage}{\linewidth}
\textbf{Step 1: Define and model.} $X$ indicates A; support $\{0,1\}$, no monetary units, parameter $p=0.3$. A and its complement are exclusive and exhaustive, so the Bernoulli model applies.
\end{minipage}\par\medskip
\par\noindent\begin{minipage}{\linewidth}
\textbf{Step 2: Write the formula.} $P(X=x)=p^x(1-p)^{1-x}$ for $x\in\{0,1\}$. The notation $x\in\{0,1\}$ means x is either 0 or 1; a nonzero number to power zero is 1.
\end{minipage}\par\medskip
\par\noindent\begin{minipage}{\linewidth}
\textbf{Step 3: Substitute.} $P(X=1)=0.3^1(0.7)^0=0.3000$, and $P(X=0)=0.3^0(0.7)^1=0.7000$.
\end{minipage}\par\medskip
\par\noindent\begin{minipage}{\linewidth}
\textbf{Step 4: Check and interpret.} $0.3+0.7=1$. The value 1 labels A; it does not measure how large A is.
\end{minipage}\par\medskip
\subsection{Bridge to the general model}
\[X\sim\operatorname{Bernoulli}(p),\qquad P(X=x)=p^x(1-p)^{1-x},\quad x\in\{0,1\}.\]
The symbol $\sim$ means has the distribution. Use this for one yes/no event during a specified period. The exponents select either $p$ or $1-p$. At $p=0$ or 1 use direct masses $P(X=1)=p$ and $P(X=0)=1-p$ to avoid $0^0$. A count from several observations needs a different model.
\Needspace{14\baselineskip}
\subsection{Advanced general problem: reversing success and measuring money}
$X$ indicates an adverse event of probability 0.2. Define $Z=1-X$ and $Y=12-50X$ USD.

\textbf{Questions/Tasks}
\begin{enumerate}
\item Find both indicator PMFs, the monetary support and $P(Y<0)$. Explain whether $Y$ is Bernoulli.
\end{enumerate}
\subsubsection*{Complete worked solution}
\par\noindent\begin{minipage}{\linewidth}
\textbf{Step 1: Define.} $X=1$ labels the adverse event; $Z=1$ labels its absence. Both have support $\{0,1\}$. $Y$ measures money, with support $\{-38,12\}$ USD.
\end{minipage}\par\medskip
\par\noindent\begin{minipage}{\linewidth}
\textbf{Step 2: Model and reverse.} $X\sim\operatorname{Bernoulli}(0.2)$, $Z\sim\operatorname{Bernoulli}(0.8)$. One event is being observed, so no trial independence assumption is needed.
\end{minipage}\par\medskip
\par\noindent\begin{minipage}{\linewidth}
\textbf{Step 3: Calculate.} $P(Z=1)=1-p=1-0.2=0.8$, $P(Z=0)=0.2$. Also $Y(1)=12-50=-38$ and $Y(0)=12$. Hence $P(Y<0)=P(X=1)=0.2000$; $P(Y=12)=0.8$.
\end{minipage}\par\medskip
\par\noindent\begin{minipage}{\linewidth}
\textbf{Step 4: Check and interpret.} Each PMF totals 1. $Y$ is a monetary transformation of an indicator, but it is not itself Bernoulli because its values are not 0 and 1. Reversing labels does not change the event causing the loss.
\end{minipage}\par\medskip
\section{Notation reference and practice procedure}
\begin{center}
\begin{tabular}{lp{0.72\linewidth}}
\toprule
Symbol & Meaning and units \\ \midrule
$X$ & Indicator: 1 for the specified event, 0 otherwise \\
$p$ & Probability of the event labelled 1; no units \\
$1-p$ & Complementary probability \\
$Z=1-X$ & Reversed indicator \\
$Y=aX+b$ & Monetary consequence if $a,b$ are monetary amounts \\
\bottomrule
\end{tabular}
\end{center}
\begin{enumerate}
\item Define the random variable in words and give its units and support.
\item List parameters and check that the model assumptions match the situation.
\item Translate ordinary language into an event, paying attention to endpoints.
\item Write the symbolic formula first, then substitute without rounding intermediate values.
\item Check bounds, total probability or a complementary calculation; interpret the answer in context.
\end{enumerate}
\Needspace{19\baselineskip}
\section{Foundational contextual practice}\label{sec:t2-c3-foundational}
\Needspace{14\baselineskip}
\subsection{One loan default}
One loan has a hypothetical probability 0.04 of default within one year. Default causes a fixed loss of 800 USD; otherwise default loss is zero.

\textbf{Questions/Tasks}
\begin{enumerate}
\item Define the indicator for the loan defaults within one year; give its PMF and the financial interpretation. 
\end{enumerate}
\subsubsection*{Complete worked solution}
\par\noindent\begin{minipage}{\linewidth}
\textbf{Step 1: Define the event.} $X=1$ if the loan defaults within one year, and $X=0$ otherwise. Support is $\{0,1\}$, without monetary units; $p=0.04$.
\end{minipage}\par\medskip
\par\noindent\begin{minipage}{\linewidth}
\textbf{Step 2: Justify the model.} This is one observation with two exclusive, exhaustive labels, so it is Bernoulli. Repeated-observation independence is not needed here.
\end{minipage}\par\medskip
\par\noindent\begin{minipage}{\linewidth}
\textbf{Step 3: Translate and calculate.} $P(X=x)=p^x(1-p)^{1-x}$ for $x=0,1$. Thus $P(X=1)=0.04$, $P(X=0)=1-0.04=0.96$. $L=800X$ USD. Thus $P(L=800)=P(X=1)=0.04$ and $P(L=0)=0.96$. Mathematical success is financially undesirable.
\end{minipage}\par\medskip
\par\noindent\begin{minipage}{\linewidth}
\textbf{Step 4: Check and interpret.} The non-negative masses sum to $ 0.04+0.96=1$. The success probability is 0.0400. The indicator records occurrence; money measures its consequence.
\end{minipage}\par\medskip
\secInicio[][sec:t2-c3-foundational]
\Needspace{14\baselineskip}
\subsection{One policy claim}
One insurance policy has probability 0.12 of at least one claim this month. The simplified contract pays 300 USD once if any claim occurs.

\textbf{Questions/Tasks}
\begin{enumerate}
\item Define the indicator for at least one claim occurs this month; give its PMF and the financial interpretation. 
\end{enumerate}
\subsubsection*{Complete worked solution}
\par\noindent\begin{minipage}{\linewidth}
\textbf{Step 1: Define the event.} $X=1$ if at least one claim occurs this month, and $X=0$ otherwise. Support is $\{0,1\}$, without monetary units; $p=0.12$.
\end{minipage}\par\medskip
\par\noindent\begin{minipage}{\linewidth}
\textbf{Step 2: Justify the model.} This is one observation with two exclusive, exhaustive labels, so it is Bernoulli. Repeated-observation independence is not needed here.
\end{minipage}\par\medskip
\par\noindent\begin{minipage}{\linewidth}
\textbf{Step 3: Translate and calculate.} $P(X=x)=p^x(1-p)^{1-x}$ for $x=0,1$. Thus $P(X=1)=0.12$, $P(X=0)=1-0.12=0.88$. $Y=300X$ USD. Hence $P(Y>0)=P(X=1)=0.12$. The indicator does not count claims.
\end{minipage}\par\medskip
\par\noindent\begin{minipage}{\linewidth}
\textbf{Step 4: Check and interpret.} The non-negative masses sum to $ 0.12+0.88=1$. The success probability is 0.1200. The indicator records occurrence; money measures its consequence.
\end{minipage}\par\medskip
\secInicio[][sec:t2-c3-foundational]
\Needspace{19\baselineskip}
\section{Intermediate contextual practice}\label{sec:t2-c3-intermediate}
\Needspace{14\baselineskip}
\subsection{Invoice payment and reversal}
One invoice is paid by its due date with probability 0.85. Late payment triggers a 6 USD fee.

\textbf{Questions/Tasks}
\begin{enumerate}
\item Define the indicator for the invoice is paid on time; give its PMF and the financial interpretation. Reverse success to lateness and find the fee probability.
\end{enumerate}
\subsubsection*{Complete worked solution}
\par\noindent\begin{minipage}{\linewidth}
\textbf{Step 1: Define the event.} $X=1$ if the invoice is paid on time, and $X=0$ otherwise. Support is $\{0,1\}$, without monetary units; $p=0.85$.
\end{minipage}\par\medskip
\par\noindent\begin{minipage}{\linewidth}
\textbf{Step 2: Justify the model.} This is one observation with two exclusive, exhaustive labels, so it is Bernoulli. Repeated-observation independence is not needed here.
\end{minipage}\par\medskip
\par\noindent\begin{minipage}{\linewidth}
\textbf{Step 3: Translate and calculate.} $P(X=x)=p^x(1-p)^{1-x}$ for $x=0,1$. Thus $P(X=1)=0.85$, $P(X=0)=1-0.85=0.15$. For lateness use $Z=1-X$, with parameter $1-p=0.15$. The fee is $L=6Z$ USD, so $P(L=6)=P(Z=1)=0.15$.
\end{minipage}\par\medskip
\par\noindent\begin{minipage}{\linewidth}
\textbf{Step 4: Check and interpret.} The non-negative masses sum to $ 0.85+0.15=1$. The success probability is 0.8500. The indicator records occurrence; money measures its consequence.
\end{minipage}\par\medskip
\secInicio[][sec:t2-c3-intermediate]
\Needspace{14\baselineskip}
\subsection{Investment gain}
One investment yields 15 USD with probability 0.6 and loses 10 USD otherwise.

\textbf{Questions/Tasks}
\begin{enumerate}
\item Define the indicator for the investment has a positive gain; give its PMF and the financial interpretation. 
\end{enumerate}
\subsubsection*{Complete worked solution}
\par\noindent\begin{minipage}{\linewidth}
\textbf{Step 1: Define the event.} $X=1$ if the investment has a positive gain, and $X=0$ otherwise. Support is $\{0,1\}$, without monetary units; $p=0.6$.
\end{minipage}\par\medskip
\par\noindent\begin{minipage}{\linewidth}
\textbf{Step 2: Justify the model.} This is one observation with two exclusive, exhaustive labels, so it is Bernoulli. Repeated-observation independence is not needed here.
\end{minipage}\par\medskip
\par\noindent\begin{minipage}{\linewidth}
\textbf{Step 3: Translate and calculate.} $P(X=x)=p^x(1-p)^{1-x}$ for $x=0,1$. Thus $P(X=1)=0.6$, $P(X=0)=1-0.6=0.4$. $G=-10+25X$ USD. Therefore $P(G=15)=0.6$, $P(G=-10)=0.4$ and $P(G<0)=0.4$. The indicator value 1 does not mean 1 USD.
\end{minipage}\par\medskip
\par\noindent\begin{minipage}{\linewidth}
\textbf{Step 4: Check and interpret.} The non-negative masses sum to $ 0.6+0.4=1$. The success probability is 0.6000. The indicator records occurrence; money measures its consequence.
\end{minipage}\par\medskip
\secInicio[][sec:t2-c3-intermediate]
\Needspace{19\baselineskip}
\section{Advanced contextual practice}\label{sec:t2-c3-advanced}
\Needspace{14\baselineskip}
\subsection{Chargeback and merchant balance}
One sale brings 30 USD initially. With probability 0.08 a chargeback removes 30 USD and adds a 5 USD fee.

\textbf{Questions/Tasks}
\begin{enumerate}
\item Define the indicator for the sale produces a chargeback; give its PMF and the financial interpretation. Find the negative-balance probability and the reversed parameter.
\end{enumerate}
\subsubsection*{Complete worked solution}
\par\noindent\begin{minipage}{\linewidth}
\textbf{Step 1: Define the event.} $X=1$ if the sale produces a chargeback, and $X=0$ otherwise. Support is $\{0,1\}$, without monetary units; $p=0.08$.
\end{minipage}\par\medskip
\par\noindent\begin{minipage}{\linewidth}
\textbf{Step 2: Justify the model.} This is one observation with two exclusive, exhaustive labels, so it is Bernoulli. Repeated-observation independence is not needed here.
\end{minipage}\par\medskip
\par\noindent\begin{minipage}{\linewidth}
\textbf{Step 3: Translate and calculate.} $P(X=x)=p^x(1-p)^{1-x}$ for $x=0,1$. Thus $P(X=1)=0.08$, $P(X=0)=1-0.08=0.92$. $Y=30-35X$ USD, with values 30 and $-5$. Thus $P(Y<0)=P(X=1)=0.08$. Reversing success to no chargeback gives parameter 0.92, but the loss probability stays 0.08.
\end{minipage}\par\medskip
\par\noindent\begin{minipage}{\linewidth}
\textbf{Step 4: Check and interpret.} The non-negative masses sum to $ 0.08+0.92=1$. The success probability is 0.0800. The indicator records occurrence; money measures its consequence.
\end{minipage}\par\medskip
\secInicio[][sec:t2-c3-advanced]
\Needspace{14\baselineskip}
\subsection{Three sales categories}
One enquiry produces no sale, a basic sale or a premium sale with probabilities $0.5,0.3,0.2$. Basic revenue is 20 USD; premium revenue is 50 USD.

\textbf{Questions/Tasks}
\begin{enumerate}
\item Explain why revenue is not Bernoulli. Give PMFs for purchase and premium indicators, and find the chance of revenue at least 20 USD.
\end{enumerate}
\subsubsection*{Complete worked solution}
\par\noindent\begin{minipage}{\linewidth}
\textbf{Step 1: Define.} $R$ is revenue in USD, support $\{0,20,50\}$. Define $X=1$ for any purchase and 0 otherwise; $Z=1$ for premium and 0 otherwise. Both indicators have support $\{0,1\}$.
\end{minipage}\par\medskip
\par\noindent\begin{minipage}{\linewidth}
\textbf{Step 2: Select.} Three revenue values exclude Bernoulli for $R$. Each indicator groups the exclusive, exhaustive outcomes into two categories, so each is Bernoulli. They concern the same enquiry and are not assumed independent.
\end{minipage}\par\medskip
\par\noindent\begin{minipage}{\linewidth}
\textbf{Step 3: Calculate.} $P(X=1)=P(\text{basic})+P(\text{premium})=0.3+0.2=0.5$; $P(X=0)=0.5$. Also $P(Z=1)=0.2$, $P(Z=0)=0.8$. Therefore $P(R\geq20)=P(X=1)=0.5000$.
\end{minipage}\par\medskip
\par\noindent\begin{minipage}{\linewidth}
\textbf{Step 4: Check and interpret.} Each PMF totals 1. Premium implies purchase, but the two indicators answer different questions. Revenue measures monetary size.
\end{minipage}\par\medskip
\secInicio[][sec:t2-c3-advanced]
\section{Closing reference: key ideas and common errors}
Define success before assigning $p$. Success may mean default or loss. Reverse the label with $1-p$. Distinguish an indicator, monetary consequence and multiple-trial count. A single Bernoulli observation does not require independence.
\secInicio
\end{document}
